\documentclass[
  doc,
  floatsintext,
  longtable,
  a4paper,
  nolmodern,
  notxfonts,
  notimes,
  12pt,
  colorlinks=true,linkcolor=blue,citecolor=blue,urlcolor=blue]{apa7}

\usepackage{amsmath}
\usepackage{amssymb}

\geometry{inner=1in, outer=1in}
\fancyhfoffset[LE,RO]{0cm}


\usepackage[bidi=default]{babel}
\babelprovide[main,import]{spanish}


% get rid of language-specific shorthands (see #6817):
\let\LanguageShortHands\languageshorthands
\def\languageshorthands#1{}

\RequirePackage{longtable}
\RequirePackage{threeparttablex}

\makeatletter
\renewcommand{\paragraph}{\@startsection{paragraph}{4}{\parindent}%
	{0\baselineskip \@plus 0.2ex \@minus 0.2ex}%
	{-.5em}%
	{\normalfont\normalsize\bfseries\typesectitle}}

\renewcommand{\subparagraph}[1]{\@startsection{subparagraph}{5}{0.5em}%
	{0\baselineskip \@plus 0.2ex \@minus 0.2ex}%
	{-\z@\relax}%
	{\normalfont\normalsize\bfseries\itshape\hspace{\parindent}{#1}\textit{\addperi}}{\relax}}
\makeatother




\usepackage{longtable, booktabs, multirow, multicol, colortbl, hhline, caption, array, float, xpatch}
\usepackage{subcaption}


\renewcommand\thesubfigure{\Alph{subfigure}}
\setcounter{topnumber}{2}
\setcounter{bottomnumber}{2}
\setcounter{totalnumber}{4}
\renewcommand{\topfraction}{0.85}
\renewcommand{\bottomfraction}{0.85}
\renewcommand{\textfraction}{0.15}
\renewcommand{\floatpagefraction}{0.7}

\usepackage{tcolorbox}
\tcbuselibrary{listings,theorems, breakable, skins}
\usepackage{fontawesome5}

\definecolor{quarto-callout-color}{HTML}{909090}
\definecolor{quarto-callout-note-color}{HTML}{0758E5}
\definecolor{quarto-callout-important-color}{HTML}{CC1914}
\definecolor{quarto-callout-warning-color}{HTML}{EB9113}
\definecolor{quarto-callout-tip-color}{HTML}{00A047}
\definecolor{quarto-callout-caution-color}{HTML}{FC5300}
\definecolor{quarto-callout-color-frame}{HTML}{ACACAC}
\definecolor{quarto-callout-note-color-frame}{HTML}{4582EC}
\definecolor{quarto-callout-important-color-frame}{HTML}{D9534F}
\definecolor{quarto-callout-warning-color-frame}{HTML}{F0AD4E}
\definecolor{quarto-callout-tip-color-frame}{HTML}{02B875}
\definecolor{quarto-callout-caution-color-frame}{HTML}{FD7E14}

%\newlength\Oldarrayrulewidth
%\newlength\Oldtabcolsep


\usepackage{hyperref}



\usepackage{color}
\usepackage{fancyvrb}
\newcommand{\VerbBar}{|}
\newcommand{\VERB}{\Verb[commandchars=\\\{\}]}
\DefineVerbatimEnvironment{Highlighting}{Verbatim}{commandchars=\\\{\}}
% Add ',fontsize=\small' for more characters per line
\usepackage{framed}
\definecolor{shadecolor}{RGB}{241,243,245}
\newenvironment{Shaded}{\begin{snugshade}}{\end{snugshade}}
\newcommand{\AlertTok}[1]{\textcolor[rgb]{0.68,0.00,0.00}{#1}}
\newcommand{\AnnotationTok}[1]{\textcolor[rgb]{0.37,0.37,0.37}{#1}}
\newcommand{\AttributeTok}[1]{\textcolor[rgb]{0.40,0.45,0.13}{#1}}
\newcommand{\BaseNTok}[1]{\textcolor[rgb]{0.68,0.00,0.00}{#1}}
\newcommand{\BuiltInTok}[1]{\textcolor[rgb]{0.00,0.23,0.31}{#1}}
\newcommand{\CharTok}[1]{\textcolor[rgb]{0.13,0.47,0.30}{#1}}
\newcommand{\CommentTok}[1]{\textcolor[rgb]{0.37,0.37,0.37}{#1}}
\newcommand{\CommentVarTok}[1]{\textcolor[rgb]{0.37,0.37,0.37}{\textit{#1}}}
\newcommand{\ConstantTok}[1]{\textcolor[rgb]{0.56,0.35,0.01}{#1}}
\newcommand{\ControlFlowTok}[1]{\textcolor[rgb]{0.00,0.23,0.31}{\textbf{#1}}}
\newcommand{\DataTypeTok}[1]{\textcolor[rgb]{0.68,0.00,0.00}{#1}}
\newcommand{\DecValTok}[1]{\textcolor[rgb]{0.68,0.00,0.00}{#1}}
\newcommand{\DocumentationTok}[1]{\textcolor[rgb]{0.37,0.37,0.37}{\textit{#1}}}
\newcommand{\ErrorTok}[1]{\textcolor[rgb]{0.68,0.00,0.00}{#1}}
\newcommand{\ExtensionTok}[1]{\textcolor[rgb]{0.00,0.23,0.31}{#1}}
\newcommand{\FloatTok}[1]{\textcolor[rgb]{0.68,0.00,0.00}{#1}}
\newcommand{\FunctionTok}[1]{\textcolor[rgb]{0.28,0.35,0.67}{#1}}
\newcommand{\ImportTok}[1]{\textcolor[rgb]{0.00,0.46,0.62}{#1}}
\newcommand{\InformationTok}[1]{\textcolor[rgb]{0.37,0.37,0.37}{#1}}
\newcommand{\KeywordTok}[1]{\textcolor[rgb]{0.00,0.23,0.31}{\textbf{#1}}}
\newcommand{\NormalTok}[1]{\textcolor[rgb]{0.00,0.23,0.31}{#1}}
\newcommand{\OperatorTok}[1]{\textcolor[rgb]{0.37,0.37,0.37}{#1}}
\newcommand{\OtherTok}[1]{\textcolor[rgb]{0.00,0.23,0.31}{#1}}
\newcommand{\PreprocessorTok}[1]{\textcolor[rgb]{0.68,0.00,0.00}{#1}}
\newcommand{\RegionMarkerTok}[1]{\textcolor[rgb]{0.00,0.23,0.31}{#1}}
\newcommand{\SpecialCharTok}[1]{\textcolor[rgb]{0.37,0.37,0.37}{#1}}
\newcommand{\SpecialStringTok}[1]{\textcolor[rgb]{0.13,0.47,0.30}{#1}}
\newcommand{\StringTok}[1]{\textcolor[rgb]{0.13,0.47,0.30}{#1}}
\newcommand{\VariableTok}[1]{\textcolor[rgb]{0.07,0.07,0.07}{#1}}
\newcommand{\VerbatimStringTok}[1]{\textcolor[rgb]{0.13,0.47,0.30}{#1}}
\newcommand{\WarningTok}[1]{\textcolor[rgb]{0.37,0.37,0.37}{\textit{#1}}}

\providecommand{\tightlist}{%
  \setlength{\itemsep}{0pt}\setlength{\parskip}{0pt}}
\usepackage{longtable,booktabs,array}
\newcounter{none} % for unnumbered tables
\usepackage{calc} % for calculating minipage widths
% Correct order of tables after \paragraph or \subparagraph
\usepackage{etoolbox}
\makeatletter
\patchcmd\longtable{\par}{\if@noskipsec\mbox{}\fi\par}{}{}
\makeatother
% Allow footnotes in longtable head/foot
\IfFileExists{footnotehyper.sty}{\usepackage{footnotehyper}}{\usepackage{footnote}}
\makesavenoteenv{longtable}

\usepackage{graphicx}
\makeatletter
\newsavebox\pandoc@box
\newcommand*\pandocbounded[1]{% scales image to fit in text height/width
  \sbox\pandoc@box{#1}%
  \Gscale@div\@tempa{\textheight}{\dimexpr\ht\pandoc@box+\dp\pandoc@box\relax}%
  \Gscale@div\@tempb{\linewidth}{\wd\pandoc@box}%
  \ifdim\@tempb\p@<\@tempa\p@\let\@tempa\@tempb\fi% select the smaller of both
  \ifdim\@tempa\p@<\p@\scalebox{\@tempa}{\usebox\pandoc@box}%
  \else\usebox{\pandoc@box}%
  \fi%
}
% Set default figure placement to htbp
\def\fps@figure{htbp}
\makeatother







\usepackage{newtx}

\defaultfontfeatures{Scale=MatchLowercase}
\defaultfontfeatures[\rmfamily]{Ligatures=TeX,Scale=1}





\title{Asignación múltiple de variables}


\shorttitle{ASIGNACIÓN MÚLTIPLE DE VARIABLES}


\usepackage{etoolbox}



\ccoppy{\textcopyright~2026}



\author{Edison Achalma}



\affiliation{
{Escuela Profesional de Economía, Universidad Nacional de San Cristóbal
de Huamanga}}




\leftheader{Achalma}

\date{2026-10-20}


\abstract{Este abstract será actualizado una vez que se complete el
contenido final del artículo. }



\authornote{\par{\addORCIDlink{Edison Achalma}{0000-0001-6996-3364}} 
\par{ }
\par{   El autor no tiene conflictos de interés que revelar.    Los
roles de autor se clasificaron utilizando la taxonomía de roles de
colaborador (CRediT; https://credit.niso.org/) de la siguiente
manera:  Edison Achalma:   conceptualización, metodología, análisis
formal, investigación, recursos, curación de
datos, redacción, visualización, supervisión, administración del
proyecto}
\par{La correspondencia relativa a este artículo debe dirigirse a Edison
Achalma, Escuela Profesional de Economía, Universidad Nacional de San
Cristóbal de Huamanga, Portal Independencia N°
57, Ayacucho, AYA 5001, Perú, Email: \href{mailto:elmer.achalma.09@unsch.edu.pe}{elmer.achalma.09@unsch.edu.pe}}
}

\makeatletter
\let\endoldlt\endlongtable
\def\endlongtable{
\hline
\endoldlt
}
\makeatother

\urlstyle{same}



\hypersetup{
  pdftitle={Asignaci\'{o}n m\'{u}ltiple de variables},
  pdfauthor={Autor}
}
\makeatletter
\@ifpackageloaded{caption}{}{\usepackage{caption}}
\AtBeginDocument{%
\ifdefined\contentsname
  \renewcommand*\contentsname{Tabla de contenidos}
\else
  \newcommand\contentsname{Tabla de contenidos}
\fi
\ifdefined\listfigurename
  \renewcommand*\listfigurename{Lista de Figuras}
\else
  \newcommand\listfigurename{Lista de Figuras}
\fi
\ifdefined\listtablename
  \renewcommand*\listtablename{Lista de Tablas}
\else
  \newcommand\listtablename{Lista de Tablas}
\fi
\ifdefined\figurename
  \renewcommand*\figurename{Figura}
\else
  \newcommand\figurename{Figura}
\fi
\ifdefined\tablename
  \renewcommand*\tablename{Tabla}
\else
  \newcommand\tablename{Tabla}
\fi
}
\@ifpackageloaded{float}{}{\usepackage{float}}
\floatstyle{ruled}
\@ifundefined{c@chapter}{\newfloat{codelisting}{h}{lop}}{\newfloat{codelisting}{h}{lop}[chapter]}
\floatname{codelisting}{Listado}
\newcommand*\listoflistings{\listof{codelisting}{Lista de Listados}}
\makeatother
\makeatletter
\makeatother
\makeatletter
\@ifpackageloaded{caption}{}{\usepackage{caption}}
\@ifpackageloaded{subcaption}{}{\usepackage{subcaption}}
\makeatother
\makeatletter
\@ifpackageloaded{fontawesome5}{}{\usepackage{fontawesome5}}
\makeatother

% From https://tex.stackexchange.com/a/645996/211326
%%% apa7 doesn't want to add appendix section titles in the toc
%%% let's make it do it
\makeatletter
\xpatchcmd{\appendix}
  {\par}
  {\addcontentsline{toc}{section}{\@currentlabelname}\par}
  {}{}
\makeatother

%% Disable longtable counter
%% https://tex.stackexchange.com/a/248395/211326

\usepackage{etoolbox}

\makeatletter
\patchcmd{\LT@caption}
  {\bgroup}
  {\bgroup\global\LTpatch@captiontrue}
  {}{}
\patchcmd{\longtable}
  {\par}
  {\par\global\LTpatch@captionfalse}
  {}{}
\apptocmd{\endlongtable}
  {\ifLTpatch@caption\else\addtocounter{table}{-1}\fi}
  {}{}
\newif\ifLTpatch@caption
\makeatother

\begin{document}

\maketitle


\hypertarget{toc}{}
\tableofcontents
\newpage
\section[Introduction]{Asignación múltiple de variables}

\setcounter{secnumdepth}{3}

\setlength\LTleft{0pt}




\subsection{🎯 Motivación}\label{motivaciuxf3n}

Python permite asignar valores a múltiples variables en una sola línea,
haciendo el código más compacto y legible. Esta característica es
especialmente útil cuando trabajas con colecciones de datos,
intercambias valores o inicializas varias variables con el mismo valor.
Dominar estas técnicas hace tu código más \textbf{Pythonico} y
eficiente.

\subsection{💡 Intuición}\label{intuiciuxf3n}

Imagina que tienes tres cajas (variables) y tres objetos (valores).
Normalmente, pondrías un objeto en cada caja por separado. La asignación
múltiple es como hacerlo todo de una vez: colocas el primer objeto en la
primera caja, el segundo en la segunda, y el tercero en la tercera,
simultáneamente. Es más rápido y más limpio.

\subsection{🧠 Modelo mental}\label{modelo-mental}

Python evalúa el lado derecho primero y luego asigna a la izquierda:

\begin{enumerate}
\def\labelenumi{\arabic{enumi}.}
\tightlist
\item
  \textbf{Asignación múltiple}: los valores se asignan en orden a las
  variables correspondientes.
\item
  \textbf{Desempaquetado}: los elementos de una colección se
  ``desempaquetan'' en variables individuales.
\item
  \textbf{Intercambio}: Python permite intercambiar valores de manera
  elegante sin variable temporal.
\end{enumerate}

\subsection{📐 Desarrollo}\label{desarrollo}

\subsubsection{1. Asignación múltiple en una
línea}\label{asignaciuxf3n-muxfaltiple-en-una-luxednea}

Asigna diferentes valores a diferentes variables:

\begin{Shaded}
\begin{Highlighting}[]
\NormalTok{x, y, z }\OperatorTok{=} \StringTok{"Manzana"}\NormalTok{, }\StringTok{"Banana"}\NormalTok{, }\StringTok{"Cereza"}
\BuiltInTok{print}\NormalTok{(x)  }\CommentTok{\# Manzana}
\BuiltInTok{print}\NormalTok{(y)  }\CommentTok{\# Banana}
\BuiltInTok{print}\NormalTok{(z)  }\CommentTok{\# Cereza}

\CommentTok{\# También funciona con números}
\NormalTok{a, b, c }\OperatorTok{=} \DecValTok{1}\NormalTok{, }\DecValTok{2}\NormalTok{, }\DecValTok{3}
\BuiltInTok{print}\NormalTok{(a, b, c)  }\CommentTok{\# 1 2 3}

\CommentTok{\# Y con tipos mixtos}
\NormalTok{nombre, edad, ciudad }\OperatorTok{=} \StringTok{"Ana"}\NormalTok{, }\DecValTok{25}\NormalTok{, }\StringTok{"Madrid"}
\end{Highlighting}
\end{Shaded}

⚠️ \textbf{Importante:} el número de variables debe coincidir con el
número de valores:

\begin{Shaded}
\begin{Highlighting}[]
\CommentTok{\# ❌ Error: demasiados valores}
\NormalTok{x, y }\OperatorTok{=} \DecValTok{1}\NormalTok{, }\DecValTok{2}\NormalTok{, }\DecValTok{3}  \CommentTok{\# ValueError: too many values to unpack}

\CommentTok{\# ❌ Error: pocos valores}
\NormalTok{x, y, z }\OperatorTok{=} \DecValTok{1}\NormalTok{, }\DecValTok{2}  \CommentTok{\# ValueError: not enough values to unpack}
\end{Highlighting}
\end{Shaded}

\subsubsection{2. Mismo valor para múltiples
variables}\label{mismo-valor-para-muxfaltiples-variables}

Asigna el mismo valor a varias variables:

\begin{Shaded}
\begin{Highlighting}[]
\NormalTok{x }\OperatorTok{=}\NormalTok{ y }\OperatorTok{=}\NormalTok{ z }\OperatorTok{=} \StringTok{"Naranja"}
\BuiltInTok{print}\NormalTok{(x)  }\CommentTok{\# Naranja}
\BuiltInTok{print}\NormalTok{(y)  }\CommentTok{\# Naranja}
\BuiltInTok{print}\NormalTok{(z)  }\CommentTok{\# Naranja}

\CommentTok{\# Con números}
\NormalTok{a }\OperatorTok{=}\NormalTok{ b }\OperatorTok{=}\NormalTok{ c }\OperatorTok{=} \DecValTok{0}
\BuiltInTok{print}\NormalTok{(a, b, c)  }\CommentTok{\# 0 0 0}
\end{Highlighting}
\end{Shaded}

\subsubsection{3. Desempaquetado
(unpacking)}\label{desempaquetado-unpacking}

Extrae valores de una colección (lista, tupla, etc.) en variables
individuales:

\begin{Shaded}
\begin{Highlighting}[]
\CommentTok{\# Desempaquetando una lista}
\NormalTok{frutas }\OperatorTok{=}\NormalTok{ [}\StringTok{"manzana"}\NormalTok{, }\StringTok{"banana"}\NormalTok{, }\StringTok{"cereza"}\NormalTok{]}
\NormalTok{x, y, z }\OperatorTok{=}\NormalTok{ frutas}
\BuiltInTok{print}\NormalTok{(x)  }\CommentTok{\# manzana}
\BuiltInTok{print}\NormalTok{(y)  }\CommentTok{\# banana}
\BuiltInTok{print}\NormalTok{(z)  }\CommentTok{\# cereza}

\CommentTok{\# Desempaquetando una tupla}
\NormalTok{coordenadas }\OperatorTok{=}\NormalTok{ (}\DecValTok{10}\NormalTok{, }\DecValTok{20}\NormalTok{)}
\NormalTok{x, y }\OperatorTok{=}\NormalTok{ coordenadas}
\BuiltInTok{print}\NormalTok{(x)  }\CommentTok{\# 10}
\BuiltInTok{print}\NormalTok{(y)  }\CommentTok{\# 20}
\end{Highlighting}
\end{Shaded}

\subsubsection{4. Intercambio de valores
(swapping)}\label{intercambio-de-valores-swapping}

Python permite intercambiar valores de forma elegante:

\begin{Shaded}
\begin{Highlighting}[]
\CommentTok{\# Método tradicional (en otros lenguajes)}
\NormalTok{a }\OperatorTok{=} \DecValTok{5}
\NormalTok{b }\OperatorTok{=} \DecValTok{10}
\NormalTok{temp }\OperatorTok{=}\NormalTok{ a}
\NormalTok{a }\OperatorTok{=}\NormalTok{ b}
\NormalTok{b }\OperatorTok{=}\NormalTok{ temp}
\BuiltInTok{print}\NormalTok{(a, b)  }\CommentTok{\# 10 5}

\CommentTok{\# Método Pythonico}
\NormalTok{a, b }\OperatorTok{=} \DecValTok{5}\NormalTok{, }\DecValTok{10}
\NormalTok{a, b }\OperatorTok{=}\NormalTok{ b, a  }\CommentTok{\# Intercambio en una línea}
\BuiltInTok{print}\NormalTok{(a, b)  }\CommentTok{\# 10 5}
\end{Highlighting}
\end{Shaded}

\subsubsection{\texorpdfstring{5. Desempaquetado con asterisco
(\texttt{*})}{5. Desempaquetado con asterisco (*)}}\label{desempaquetado-con-asterisco}

Recoge múltiples valores en una lista:

\begin{Shaded}
\begin{Highlighting}[]
\CommentTok{\# Recoge el resto de elementos}
\NormalTok{a, b, }\OperatorTok{*}\NormalTok{resto }\OperatorTok{=} \DecValTok{1}\NormalTok{, }\DecValTok{2}\NormalTok{, }\DecValTok{3}\NormalTok{, }\DecValTok{4}\NormalTok{, }\DecValTok{5}
\BuiltInTok{print}\NormalTok{(a)      }\CommentTok{\# 1}
\BuiltInTok{print}\NormalTok{(b)      }\CommentTok{\# 2}
\BuiltInTok{print}\NormalTok{(resto)  }\CommentTok{\# [3, 4, 5]}

\CommentTok{\# También funciona en medio}
\NormalTok{a, }\OperatorTok{*}\NormalTok{medio, b }\OperatorTok{=} \DecValTok{1}\NormalTok{, }\DecValTok{2}\NormalTok{, }\DecValTok{3}\NormalTok{, }\DecValTok{4}\NormalTok{, }\DecValTok{5}
\BuiltInTok{print}\NormalTok{(a)      }\CommentTok{\# 1}
\BuiltInTok{print}\NormalTok{(medio)  }\CommentTok{\# [2, 3, 4]}
\BuiltInTok{print}\NormalTok{(b)      }\CommentTok{\# 5}
\end{Highlighting}
\end{Shaded}

\subsection{🔬 Ejemplos}\label{ejemplos}

\textbf{Mínimo ---} Asignación básica:

\begin{Shaded}
\begin{Highlighting}[]
\CommentTok{\# Múltiples valores}
\NormalTok{nombre, edad, ciudad }\OperatorTok{=} \StringTok{"Ana"}\NormalTok{, }\DecValTok{25}\NormalTok{, }\StringTok{"Madrid"}
\BuiltInTok{print}\NormalTok{(nombre, edad, ciudad)  }\CommentTok{\# Ana 25 Madrid}

\CommentTok{\# Mismo valor}
\NormalTok{x }\OperatorTok{=}\NormalTok{ y }\OperatorTok{=}\NormalTok{ z }\OperatorTok{=} \DecValTok{100}
\BuiltInTok{print}\NormalTok{(x, y, z)  }\CommentTok{\# 100 100 100}
\end{Highlighting}
\end{Shaded}

\textbf{Progresivo ---} Desempaquetado:

\begin{Shaded}
\begin{Highlighting}[]
\CommentTok{\# Desempaquetado de lista}
\NormalTok{colores }\OperatorTok{=}\NormalTok{ [}\StringTok{"rojo"}\NormalTok{, }\StringTok{"verde"}\NormalTok{, }\StringTok{"azul"}\NormalTok{]}
\NormalTok{primero, segundo, tercero }\OperatorTok{=}\NormalTok{ colores}
\BuiltInTok{print}\NormalTok{(primero)   }\CommentTok{\# rojo}
\BuiltInTok{print}\NormalTok{(segundo)   }\CommentTok{\# verde}
\BuiltInTok{print}\NormalTok{(tercero)   }\CommentTok{\# azul}
\end{Highlighting}
\end{Shaded}

\textbf{Real ---} Sistema de ventas:

\begin{Shaded}
\begin{Highlighting}[]
\CommentTok{\# Datos de un producto}
\NormalTok{producto }\OperatorTok{=}\NormalTok{ [}\StringTok{"Laptop"}\NormalTok{, }\FloatTok{850.50}\NormalTok{, }\DecValTok{15}\NormalTok{]  }\CommentTok{\# nombre, precio, stock}
\NormalTok{nombre, precio, stock }\OperatorTok{=}\NormalTok{ producto}

\CommentTok{\# Cálculo de venta}
\NormalTok{cantidad }\OperatorTok{=} \DecValTok{3}
\NormalTok{total }\OperatorTok{=}\NormalTok{ precio }\OperatorTok{*}\NormalTok{ cantidad}
\NormalTok{nuevo\_stock }\OperatorTok{=}\NormalTok{ stock }\OperatorTok{{-}}\NormalTok{ cantidad}

\BuiltInTok{print}\NormalTok{(}\StringTok{"Producto:"}\NormalTok{, nombre)}
\BuiltInTok{print}\NormalTok{(}\StringTok{"Precio unitario: $"}\NormalTok{, precio)}
\BuiltInTok{print}\NormalTok{(}\StringTok{"Total venta: $"}\NormalTok{, total)}
\BuiltInTok{print}\NormalTok{(}\StringTok{"Stock restante:"}\NormalTok{, nuevo\_stock)}
\end{Highlighting}
\end{Shaded}

\subsection{⚠️ Errores comunes}\label{errores-comunes}

{\def\LTcaptype{none} % do not increment counter
\begin{longtable}[]{@{}
  >{\raggedright\arraybackslash}p{(\linewidth - 4\tabcolsep) * \real{0.2409}}
  >{\raggedright\arraybackslash}p{(\linewidth - 4\tabcolsep) * \real{0.3650}}
  >{\raggedright\arraybackslash}p{(\linewidth - 4\tabcolsep) * \real{0.3942}}@{}}
\toprule\noalign{}
\begin{minipage}[b]{\linewidth}\raggedright
Error
\end{minipage} & \begin{minipage}[b]{\linewidth}\raggedright
Síntoma
\end{minipage} & \begin{minipage}[b]{\linewidth}\raggedright
Cómo evitarlo
\end{minipage} \\
\midrule\noalign{}
\endhead
\bottomrule\noalign{}
\endlastfoot
\textbf{Número de variables ≠ valores} &
\texttt{ValueError:\ too\ many/few\ values\ to\ unpack} & Cuenta
variables y valores \\
\textbf{Desempaquetar no iterable} &
\texttt{TypeError:\ cannot\ unpack\ non-iterable\ int\ object} & Solo
puedes desempaquetar colecciones \\
\textbf{Usar \texttt{*} sin necesidad} & Funciona, pero puede ser
confuso & Usa \texttt{*} solo cuando el número de elementos es
variable \\
\textbf{Olvidar el orden} & Asignación incorrecta & Los valores se
asignan en orden \\
\end{longtable}
}

\subsection{✅ Buenas prácticas}\label{buenas-pruxe1cticas}

\begin{enumerate}
\def\labelenumi{\arabic{enumi}.}
\item
  \textbf{Usa asignación múltiple para intercambiar valores}:

\begin{Shaded}
\begin{Highlighting}[]
\CommentTok{\# ✅ Más limpio}
\NormalTok{a, b }\OperatorTok{=}\NormalTok{ b, a}

\CommentTok{\# ❌ Menos elegante}
\NormalTok{temp }\OperatorTok{=}\NormalTok{ a}
\NormalTok{a }\OperatorTok{=}\NormalTok{ b}
\NormalTok{b }\OperatorTok{=}\NormalTok{ temp}
\end{Highlighting}
\end{Shaded}
\item
  \textbf{Desempaqueta colecciones para mejorar legibilidad}:

\begin{Shaded}
\begin{Highlighting}[]
\CommentTok{\# ✅ Legible}
\NormalTok{coordenadas }\OperatorTok{=}\NormalTok{ (}\DecValTok{10}\NormalTok{, }\DecValTok{20}\NormalTok{, }\DecValTok{30}\NormalTok{)}
\NormalTok{x, y, z }\OperatorTok{=}\NormalTok{ coordenadas}

\CommentTok{\# ❌ Confuso}
\NormalTok{x }\OperatorTok{=}\NormalTok{ coordenadas[}\DecValTok{0}\NormalTok{]}
\NormalTok{y }\OperatorTok{=}\NormalTok{ coordenadas[}\DecValTok{1}\NormalTok{]}
\NormalTok{z }\OperatorTok{=}\NormalTok{ coordenadas[}\DecValTok{2}\NormalTok{]}
\end{Highlighting}
\end{Shaded}
\item
  \textbf{Usa \texttt{\_} para valores que no necesitas}:

\begin{Shaded}
\begin{Highlighting}[]
\CommentTok{\# Ignora el segundo valor}
\NormalTok{x, \_, z }\OperatorTok{=} \DecValTok{1}\NormalTok{, }\DecValTok{2}\NormalTok{, }\DecValTok{3}
\BuiltInTok{print}\NormalTok{(x, z)  }\CommentTok{\# 1 3}
\end{Highlighting}
\end{Shaded}
\item
  \textbf{No abuses de la asignación múltiple}: si tienes muchas
  variables (\textgreater3), considera usar una estructura de datos.
\end{enumerate}

\subsection{📝 Ejercicio}\label{ejercicio}

\textbf{Básico} --- Asignación y desempaquetado

Escribe un programa que:

\begin{enumerate}
\def\labelenumi{\arabic{enumi}.}
\tightlist
\item
  Cree una lista \texttt{datos\_usuario} con: nombre, edad, email,
  ciudad
\item
  Desempaquete la lista en variables individuales
\item
  Cree una variable \texttt{telefono} y asígnale el valor ``555-1234''
\item
  Cree tres variables \texttt{a}, \texttt{b}, \texttt{c} con el mismo
  valor 0
\item
  Intercambie los valores de \texttt{x\ =\ 10} y \texttt{y\ =\ 20} en
  una línea
\item
  Imprima todos los resultados
\end{enumerate}

\begin{center}\rule{0.5\linewidth}{0.5pt}\end{center}

\textbf{Solución (no mires hasta haberlo intentado):}

\begin{Shaded}
\begin{Highlighting}[]
\CommentTok{\# 1. Lista de datos}
\NormalTok{datos\_usuario }\OperatorTok{=}\NormalTok{ [}\StringTok{"Ana Martínez"}\NormalTok{, }\DecValTok{28}\NormalTok{, }\StringTok{"ana@email.com"}\NormalTok{, }\StringTok{"Barcelona"}\NormalTok{]}

\CommentTok{\# 2. Desempaquetado}
\NormalTok{nombre, edad, email, ciudad }\OperatorTok{=}\NormalTok{ datos\_usuario}

\CommentTok{\# 3. Variable individual}
\NormalTok{telefono }\OperatorTok{=} \StringTok{"555{-}1234"}

\CommentTok{\# 4. Mismo valor para múltiples variables}
\NormalTok{a }\OperatorTok{=}\NormalTok{ b }\OperatorTok{=}\NormalTok{ c }\OperatorTok{=} \DecValTok{0}

\CommentTok{\# 5. Intercambio de valores}
\NormalTok{x, y }\OperatorTok{=} \DecValTok{10}\NormalTok{, }\DecValTok{20}
\NormalTok{x, y }\OperatorTok{=}\NormalTok{ y, x}

\CommentTok{\# 6. Resultados}
\BuiltInTok{print}\NormalTok{(}\StringTok{"Nombre:"}\NormalTok{, nombre)}
\BuiltInTok{print}\NormalTok{(}\StringTok{"Edad:"}\NormalTok{, edad)}
\BuiltInTok{print}\NormalTok{(}\StringTok{"Email:"}\NormalTok{, email)}
\BuiltInTok{print}\NormalTok{(}\StringTok{"Ciudad:"}\NormalTok{, ciudad)}
\BuiltInTok{print}\NormalTok{(}\StringTok{"Teléfono:"}\NormalTok{, telefono)}
\BuiltInTok{print}\NormalTok{(}\StringTok{"a, b, c:"}\NormalTok{, a, b, c)}
\BuiltInTok{print}\NormalTok{(}\StringTok{"x, y después del intercambio:"}\NormalTok{, x, y)}
\end{Highlighting}
\end{Shaded}

\subsection{🗂️ Síntesis}\label{suxedntesis}

\begin{itemize}
\tightlist
\item
  \textbf{Asignación múltiple}: \texttt{x,\ y,\ z\ =\ 1,\ 2,\ 3}
  (coincidir número de variables y valores).
\item
  \textbf{Mismo valor}: \texttt{x\ =\ y\ =\ z\ =\ 0} (todas apuntan al
  mismo valor).
\item
  \textbf{Desempaquetado}: extrae elementos de listas/tuplas en
  variables.
\item
  \textbf{Intercambio}: \texttt{a,\ b\ =\ b,\ a} para intercambiar
  valores (muy Pythonico).
\item
  \textbf{Asterisco (\texttt{*})}: recolecta múltiples valores en una
  lista.
\item
  \textbf{Guión bajo (\texttt{\_})}: ignora valores no deseados.
\end{itemize}

\subsection{🔗 Qué sigue}\label{quuxe9-sigue}

\begin{itemize}
\tightlist
\item
  \textbf{Relacionado:} listas y tuplas (estructuras de datos),
  desempaquetado avanzado
\item
  \textbf{Siguiente tema:} Tipos de datos --- aprenderás a manejar
  diferentes tipos de datos como strings, números y booleanos.
\end{itemize}

\section{Publicaciones Similares}\label{publicaciones-similares}

Si te interesó este artículo, te recomendamos que explores otros blogs y
recursos relacionados que pueden ampliar tus conocimientos. Aquí te dejo
algunas sugerencias:

\begin{enumerate}
\def\labelenumi{\arabic{enumi}.}
\tightlist
\item
  \href{https://numerus-scriptum.netlify.app/python/2020-06-19-instalacion-de-anaconda/index.pdf}{\faIcon{file-pdf}}
  \href{https://numerus-scriptum.netlify.app/python/2020-06-19-instalacion-de-anaconda}{Instalacion
  De Anaconda}
\item
  \href{https://numerus-scriptum.netlify.app/python/2020-06-20-configurar-entorno-virtual-python-anaconda/index.pdf}{\faIcon{file-pdf}}
  \href{https://numerus-scriptum.netlify.app/python/2020-06-20-configurar-entorno-virtual-python-anaconda}{Configurar
  Entorno Virtual Python Anaconda}
\item
  \href{https://numerus-scriptum.netlify.app/python/2021-04-17-01-introducion-a-la-programacion-con-python/index.pdf}{\faIcon{file-pdf}}
  \href{https://numerus-scriptum.netlify.app/python/2021-04-17-01-introducion-a-la-programacion-con-python}{01
  Introducion A La Programacion Con Python}
\item
  \href{https://numerus-scriptum.netlify.app/python/2021-05-31-02-variables-expresiones-y-statements-con-python/index.pdf}{\faIcon{file-pdf}}
  \href{https://numerus-scriptum.netlify.app/python/2021-05-31-02-variables-expresiones-y-statements-con-python}{02
  Variables Expresiones Y Statements Con Python}
\item
  \href{https://numerus-scriptum.netlify.app/python/2021-06-07-03-objetos-de-python/index.pdf}{\faIcon{file-pdf}}
  \href{https://numerus-scriptum.netlify.app/python/2021-06-07-03-objetos-de-python}{03
  Objetos De Python}
\item
  \href{https://numerus-scriptum.netlify.app/python/2021-06-14-04-ejecucion-condicional-con-python/index.pdf}{\faIcon{file-pdf}}
  \href{https://numerus-scriptum.netlify.app/python/2021-06-14-04-ejecucion-condicional-con-python}{04
  Ejecucion Condicional Con Python}
\item
  \href{https://numerus-scriptum.netlify.app/python/2021-06-21-05-iteraciones-con-python/index.pdf}{\faIcon{file-pdf}}
  \href{https://numerus-scriptum.netlify.app/python/2021-06-21-05-iteraciones-con-python}{05
  Iteraciones Con Python}
\item
  \href{https://numerus-scriptum.netlify.app/python/2021-08-16-06-funciones-con-python/index.pdf}{\faIcon{file-pdf}}
  \href{https://numerus-scriptum.netlify.app/python/2021-08-16-06-funciones-con-python}{06
  Funciones Con Python}
\item
  \href{https://numerus-scriptum.netlify.app/python/2021-08-23-07-dataframes-con-python/index.pdf}{\faIcon{file-pdf}}
  \href{https://numerus-scriptum.netlify.app/python/2021-08-23-07-dataframes-con-python}{07
  Dataframes Con Python}
\item
  \href{https://numerus-scriptum.netlify.app/python/2021-11-29-08-prediccion-y-metrica-de-performance-con-python/index.pdf}{\faIcon{file-pdf}}
  \href{https://numerus-scriptum.netlify.app/python/2021-11-29-08-prediccion-y-metrica-de-performance-con-python}{08
  Prediccion Y Metrica De Performance Con Python}
\item
  \href{https://numerus-scriptum.netlify.app/python/2021-12-06-09-metodos-de-machine-learning-para-clasificacion-con-python/index.pdf}{\faIcon{file-pdf}}
  \href{https://numerus-scriptum.netlify.app/python/2021-12-06-09-metodos-de-machine-learning-para-clasificacion-con-python}{09
  Metodos De Machine Learning Para Clasificacion Con Python}
\item
  \href{https://numerus-scriptum.netlify.app/python/2021-12-13-10-metodos-de-machine-learning-para-regresion-con-python/index.pdf}{\faIcon{file-pdf}}
  \href{https://numerus-scriptum.netlify.app/python/2021-12-13-10-metodos-de-machine-learning-para-regresion-con-python}{10
  Metodos De Machine Learning Para Regresion Con Python}
\item
  \href{https://numerus-scriptum.netlify.app/python/2025-05-10-visualizacion-de-datos-con-python/index.pdf}{\faIcon{file-pdf}}
  \href{https://numerus-scriptum.netlify.app/python/2025-05-10-visualizacion-de-datos-con-python}{Visualizacion
  De Datos Con Python}
\end{enumerate}

Esperamos que encuentres estas publicaciones igualmente interesantes y
útiles. ¡Disfruta de la lectura!






\end{document}
